\subsection{Vòng đời ứng dụng và quản lý trạng thái}
\label{sec:lifecycle-state}

\subsubsection{Luồng thực thi (Execution Flow)}
\label{sec:execution-flow}

Luồng khởi động của OneSpace:

\begin{enumerate}
\item \textbf{Launch app}: Android khởi chạy MainActivity, OneSpaceApplication khởi tạo Hilt dependencies.
\item \textbf{Load preferences and saved session}: đọc DataStore (cài đặt người dùng) và Room/DataStore (phiên hoạt động trước đó).
\item \textbf{Kiểm tra phiên đăng nhập}: nếu có phiên hợp lệ $\rightarrow$ vào Home trực tiếp; nếu không $\rightarrow$ vào Guest Mode hoặc màn Đăng nhập.
\item \textbf{Home Screen}: hiển thị nội dung gợi ý theo sở thích; người dùng chọn phân hệ (Video, Music, News, Games).
\item \textbf{SessionCoordinator} lưu hoặc khôi phục trạng thái phân hệ tương ứng.
\end{enumerate}

Nhóm khuyên \textbf{Guest Mode} trong MVP (tương ứng F-COM-18) để người dùng trải nghiệm đủ bốn phân hệ mà không bị phụ thuộc backend. Authentication là module mở rộng nếu assignment không yêu cầu tài khoản bắt buộc.

\subsubsection{Quy tắc giữ trạng thái}
\label{sec:state-rules}

Điểm rất quan trọng: \textbf{chuyển tab trong app khác với đưa toàn bộ ứng dụng xuống background}. Chuyển \texttt{Video $\rightarrow$ Tin} vẫn nằm trong cùng Activity, nên Floating Video tiếp tục phát. Nhưng bấm Home đưa Activity xuống nền thì MVP \textbf{pause video} để tiết kiệm pin, trừ khi người dùng kích hoạt chế độ riêng.

Bảng~\ref{tab:state-rules} định nghĩa trạng thái từng loại nội dung trong ba tình huống:

\begin{longtable}{>{\raggedright\arraybackslash}p{1.6cm} >{\raggedright\arraybackslash}p{4.6cm} >{\raggedright\arraybackslash}p{5.2cm} >{\raggedright\arraybackslash}p{4.6cm}}
\caption{Quy tắc giữ trạng thái theo ba tình huống}
\label{tab:state-rules}\\
\toprule
\textbf{Nội dung} & \textbf{Chuyển tab trong app} & \textbf{App vào background} & \textbf{Process bị kill} \\
\midrule
\endfirsthead
\toprule
\textbf{Nội dung} & \textbf{Chuyển tab} & \textbf{Background} & \textbf{Process kill} \\
\midrule
\endhead
\bottomrule
\endfoot
Video & Floating hoặc pause theo chính sách & Pause (trừ khi bật PiP) & Khôi phục video và vị trí, không tự phát \\
Music & Tiếp tục phát nếu không xung đột & Tiếp tục qua MediaSessionService & Lưu bài, vị trí và queue; khôi phục khi mở lại \\
News & Giữ vị trí cuộn & Giữ trạng thái đọc & Phục hồi article ID và vị trí đọc \\
Games & Pause + save & Pause + save & Phục hồi từ bản lưu gần nhất \\
\end{longtable}

\subsubsection{State nào được lưu ở đâu}
\label{sec:state-storage}

\begin{longtable}{>{\raggedright\arraybackslash}p{5.2cm} >{\raggedright\arraybackslash}p{9.8cm}}
\caption{Loại trạng thái và cách lưu}
\label{tab:state-storage}\\
\toprule
\textbf{Loại trạng thái} & \textbf{Cách lưu} \\
\midrule
\endfirsthead
\toprule
\textbf{Loại trạng thái} & \textbf{Cách lưu} \\
\midrule
\endhead
\bottomrule
\endfoot
Nút đang chọn, popup đang mở & Compose \texttt{rememberSaveable} \\
UI State màn hình & ViewModel + StateFlow \\
Navigation arguments & Navigation / SavedStateHandle \\
Vị trí video, bài nhạc, queue & Repository + Room/DataStore \\
Nội dung tin đã cache & Room \\
Game board, score, level & Room \\
Cài đặt mute, chất lượng media & DataStore \\
Player còn hoạt động & Playback controller / service~--~\textbf{không} lưu đối tượng Player xuống database \\
\end{longtable}

Lưu ý quan trọng: ViewModel tồn tại qua configuration changes thông thường (xoay màn hình), nhưng \textbf{không tồn tại qua process death}. Vì vậy game state và tiến độ media không thể chỉ giữ trong ViewModel~--~phải lưu xuống Room hoặc DataStore. Cũng không nên trông chờ \texttt{onDestroy()} được gọi trước khi hệ điều hành kill process; nên lưu các mốc quan trọng ngay khi trạng thái thay đổi, thay vì chỉ lưu tất cả khi app thoát.

\subsubsection{Sơ đồ vòng đời Activity}
\label{sec:activity-lifecycle}

Vòng đời Activity chuẩn của Android:

\begin{quote}
\texttt{Created} $\rightarrow$ \texttt{Started} $\rightarrow$ \texttt{Resumed} $\rightarrow$ \texttt{Paused} $\rightarrow$ \texttt{Stopped} $\rightarrow$ \texttt{Destroyed}
\end{quote}

Sự kiện chính và ứng xử:

\begin{itemize}
\item \textbf{App loses focus} (notification kéo xuống, dialog): Activity vào \texttt{Paused}~--~video pause, nhạc vẫn phát.
\item \textbf{No longer visible} (bấm Home, chuyển app): Activity vào \texttt{Stopped}~--~video và game pause + save; nhạc tiếp tục qua MediaSessionService.
\item \textbf{System reclaims memory}: process bị kill~--~khôi phục từ Room/DataStore khi mở lại.
\item \textbf{Activity finishes}: người dùng chủ động thoát hoặc hệ thống.
\end{itemize}

Các trạng thái này chỉ đại diện vòng đời \textbf{Activity}. Phát nhạc qua MediaSessionService có vòng đời riêng, không gắn cứng với MainActivity~--~khi Activity bị destroy, service vẫn chạy.

\subsubsection{Background Audio Playback}
\label{sec:background-audio}

Phân hệ Music dùng \texttt{MediaSessionService} để duy trì media session ngoài vòng đời Activity, cung cấp điều khiển từ notification, Bluetooth và màn hình khoá. Bảng~\ref{tab:audio-rules} định nghĩa hành vi mong muốn theo từng sự kiện:

\begin{longtable}{>{\raggedright\arraybackslash}p{5.4cm} >{\raggedright\arraybackslash}p{9.6cm}}
\caption{Quy tắc audio focus theo sự kiện}
\label{tab:audio-rules}\\
\toprule
\textbf{Sự kiện} & \textbf{Hành vi mong muốn} \\
\midrule
\endfirsthead
\toprule
\textbf{Sự kiện} & \textbf{Hành vi} \\
\midrule
\endhead
\bottomrule
\endfoot
Người dùng khoá màn hình & Music tiếp tục phát \\
Bấm Home & Music tiếp tục phát \\
Cuộc gọi đến & Pause theo audio focus change \\
Rút tai nghe & Pause~--~tránh phát loa ngoài bất ngờ \\
Notification Pause & Music tạm dừng \\
Notification Next & Chuyển bài tiếp theo \\
Mở Video có âm thanh & Music pause (case C1) \\
Mở game đã mute & Music tiếp tục (case C5) \\
Người dùng chủ động Stop & Dừng session nền theo chính sách \\
\end{longtable}

Khi triển khai Media3, nhóm bật xử lý Android Audio Focus qua \texttt{setAudioAttributes(..., true)}, đồng thời dùng chính sách nội bộ (bảng~\ref{tab:audio-rules}) để phối hợp các nguồn phát trong cùng ứng dụng.

\subsubsection{Ma trận chuyển phân hệ~--~đối chiếu với case C1--C9}
\label{sec:multitasking-matrix}

Bảng~\ref{tab:multitasking-matrix} là phiên bản kỹ thuật của ma trận case C1--C9 đã trình bày ở mục~\ref{sec:mvp-scope}, gắn trực tiếp với các thành phần điều phối kiến trúc:

\begin{longtable}{>{\raggedright\arraybackslash}p{3.4cm} >{\raggedright\arraybackslash}p{4.8cm} >{\raggedright\arraybackslash}p{6.8cm}}
\caption{Ma trận chuyển phân hệ: trạng thái player/game và UI}
\label{tab:multitasking-matrix}\\
\toprule
\textbf{Chuyển từ $\rightarrow$ tới} & \textbf{Player / Game State} & \textbf{UI hiển thị} \\
\midrule
\endfirsthead
\toprule
\textbf{Chuyển từ $\rightarrow$ tới} & \textbf{Player / Game} & \textbf{UI} \\
\midrule
\endhead
\bottomrule
\endfoot
Video $\rightarrow$ Tin & Video tiếp tục (case C3) & Floating Video \\
Video $\rightarrow$ Nhạc & Video pause, Nhạc play (case C1) & Mini Music Player \\
Video $\rightarrow$ Games & Video tiếp tục nếu game mute (case C4) & Floating Video + game \\
Nhạc $\rightarrow$ Tin & Nhạc tiếp tục (case C2) & Mini Music Player \\
Nhạc $\rightarrow$ Video & Nhạc pause khi video phát tiếng & Video Player \\
Nhạc $\rightarrow$ Games & Nhạc tiếp tục nếu game mute (case C5) & Game fullscreen, nhạc nền \\
Games $\rightarrow$ Tin & Game pause + save & News Screen \\
Games $\rightarrow$ Video & Game pause + save & Video Screen \\
Tin $\rightarrow$ Games & Lưu vị trí đọc & Game Screen \\
Any $\rightarrow$ Background & Pause video/game; Nhạc tiếp tục & Android background \\
\end{longtable}

Đây là \textbf{chính sách sản phẩm} nhóm chủ động lập trình, không phải hành vi Android tự động thực hiện. Mỗi chuyển tab phải được code: pause, resume, đổi player, lưu session.

\subsubsection{Quyết định kiến trúc quan trọng}
\label{sec:adr}

Bảng~\ref{tab:adr} tổng hợp các Architecture Decision Records (ADR) nhóm cần thống nhất trước khi code, để tránh các thành viên code theo cách khác nhau:

\begin{longtable}{>{\raggedright\arraybackslash}p{1.6cm} >{\raggedright\arraybackslash}p{13.4cm}}
\caption{Architecture Decision Records}
\label{tab:adr}\\
\toprule
\textbf{ADR} & \textbf{Quyết định} \\
\midrule
\endfirsthead
\toprule
\textbf{ADR} & \textbf{Quyết định} \\
\midrule
\endhead
\bottomrule
\endfoot
ADR-01 & Android Native, Kotlin + Jetpack Compose \\
ADR-02 & Android 8.0+ (API 26+), mục tiêu RAM tối thiểu 3~GB \\
ADR-03 & MVVM + 3-layer architecture + UDF (Unidirectional Data Flow) \\
ADR-04 & Feature-based modularization \\
ADR-05 & Media3 cho audio/video playback \\
ADR-06 & MediaSessionService cho Music background \\
ADR-07 & Một nguồn audio chính có tiếng tại một thời điểm; cho phép game mute SFX \\
ADR-08 & In-app Floating Player cho Video $\rightarrow$ Tin/Games \\
ADR-09 & Room cho persisted session/game state; DataStore cho preferences \\
ADR-10 & News cache và offline-first minigame \\
ADR-11 & SessionCoordinator xử lý chuyển đổi phân hệ \\
ADR-12 & Video và game mặc định pause khi app vào background; Music được tiếp tục \\
\end{longtable}

\subsubsection{Test case kiểm chứng kiến trúc}
\label{sec:arch-test-cases}

Để chứng minh kiến trúc hoạt động đúng, nhóm đưa vào báo cáo các test case thực tế sau:

\begin{longtable}{>{\raggedright\arraybackslash}p{1.4cm} >{\raggedright\arraybackslash}p{6.4cm} >{\raggedright\arraybackslash}p{7.2cm}}
\caption{Test case kiểm chứng hành vi đa nhiệm và vòng đời}
\label{tab:arch-test-cases}\\
\toprule
\textbf{Test} & \textbf{Tình huống} & \textbf{Kết quả mong đợi} \\
\midrule
\endfirsthead
\toprule
\textbf{Test} & \textbf{Tình huống} & \textbf{Kết quả} \\
\midrule
\endhead
\bottomrule
\endfoot
TC01 & Video $\rightarrow$ Tin & Video thu nhỏ, tiếp tục phát \\
TC02 & Video $\rightarrow$ Nhạc & Video pause, Nhạc bắt đầu \\
TC03 & Nhạc $\rightarrow$ Home $\rightarrow$ khoá màn hình & Nhạc tiếp tục, điều khiển từ lock screen \\
TC04 & Nhạc $\rightarrow$ Games (game mute) & Nhạc tiếp tục \\
TC05 & Game $\rightarrow$ Tin $\rightarrow$ Game & Khôi phục đúng board và score \\
TC06 & Mất mạng khi đọc Tin & Hiển thị bài đã cache \\
TC07 & Mất mạng khi stream Video & Phát buffer còn lại, sau đó trạng thái chờ/retry \\
TC08 & Có mạng trở lại & Cho phép retry hoặc resume \\
TC09 & Activity recreate (xoay màn hình) & Trạng thái cần thiết được giữ \\
TC10 & OS kill app process & Khôi phục phiên từ local storage \\
TC11 & Có cuộc gọi khi nghe Nhạc & Audio pause/duck theo chính sách \\
TC12 & Rút tai nghe & Nhạc pause \\
TC13 & Thiết bị RAM 3~GB & Không crash/OOM trong luồng demo (Quốc Huy) \\
TC14 & Bật Battery Saver & Hạn chế tác vụ không thiết yếu \\
\end{longtable}

Để kiểm chứng các tuyên bố về pin và hiệu năng, nhóm có thể dùng Android Studio Profiler, Memory Profiler và Perfetto~--~đo khi chỉ xem video so với vừa xem floating video vừa đọc tin hoặc chơi game. Các số liệu này bổ sung cho bảng chỉ số ở mục~\ref{sec:quality-pillars}.
